% Ad Familiares 13.55 --- headnote
% ad-familiares-13-55, undated (placeholder 54 BC), Cilicia / Asia province

\textit{Cicero to Q. Minucius Thermus, propraetor of Asia,
undated --- one of the cluster of recommendations Cicero
addressed to Thermus during his own proconsulship of Cilicia
in 51--50 BC. Perseus's tradition places the letter just
before \textit{Fam.} 13.53 (\textit{paulo ante ep.\ liii}),
during the same campaign year. Cicero had met Thermus at
Ephesus and discussed the case there in person; the letter
follows up in writing what was already agreed face to face.}

\textit{The man being commended, Marcus Anneius, was Cicero's
own legate on the Cilician staff: a man Cicero had taken on
unsolicited, having refused others, and whom service in the
Parthian-frontier campaign had convinced him to rate above
all his colleagues. The dispute with the people of Sardis is
not further specified --- a civil suit of some weight, on
which Cicero has briefed Thermus orally. The shape of the
letter is the standard \textit{commendatio} doubled: a
commendation of the man, and a commendation of his cause. The
final formula --- that Anneius will not doubt your friendship
and is to be made still more your friend by this letter ---
is one of the recurring closing turns of the recommendation
genre.}
